%%%%%%%%%%%%%%%%%%%%%%%%%%%%%%%%%%%%%%%%%%%%%%%%%%%%%%%%%%%%%%
%%%%%%%%%%%%%%%%%% MA-0528 Sistemas dinámicos %%%%%%%%%%%%%%%%%%%%%%%%%%%%%%
%%%%%%%%%%%%%%%%%%%%%%%%%%%%%%%%%%%%%%%%%%%%%%%%%%%%%%%%%%%%%%
\newpage
\subsection*{MA-0528 Sistemas dinámicos}
\addcontentsline{toc}{subsection}{MA-0528 Sistemas dinámicos}

\hypertarget{MA0528}{}
\begin{refsection}
\begin{table}[H]
\centering{
\begin{tabular}{|ll|}
\hline
\textbf{Año:} IV  & \textbf{Requisitos:} MA-0515, MA-0551 y MA-0635 \\ 
\textbf{Ciclo:} Optativa  & \textbf{Correquisitos:} Ninguno \\ 
\textbf{Tipo de curso:} Teórico & \textbf{Horas presenciales por semana:} 5 \\ 
\textbf{Créditos:} 5 & \textbf{Horas de trabajo independiente por semana:} 10 \\
\textbf{Virtualidad:} P  &  \\ \hline
\end{tabular}
}
\end{table}

\subsubsection*{Descripción del curso}

Este curso ofrece una introducción rigurosa a los sistemas dinámicos y a la teoría ergódica, con énfasis en sistemas unidimensionales. Se estudian sistemas en tiempo discreto y continuo, puntos fijos y periódicos, conjuntos límite, recurrencia y conjugaciones; se analizan los difeomorfismos del intervalo, la linealización local de contracciones y la codificación simbólica de conjuntos de Cantor y de transformaciones expansoras por ramas completas.

La segunda parte desarrolla las nociones de medida invariante, recurrencia de Poincaré, promedios de Birkhoff y ergodicidad. Estas herramientas se aplican a mapas afines y expansores por ramas, al mapa de Gauss y a la familia cuadrática.

\subsubsection*{Objetivos}

\textbf{General:} Que la persona estudiante adquiera las herramientas fundamentales para analizar sistemas dinámicos desde los puntos de vista topológico, diferenciable y ergódico, reconocer estructuras recurrentes, construir modelos simbólicos, estudiar medidas invariantes e interpretar el comportamiento estadístico de las órbitas mediante promedios temporales.

\textbf{Específicos:} Al finalizar el curso se espera que la persona estudiante sea capaz de:

\begin{enumerate}
\item Formular sistemas dinámicos en tiempo discreto y continuo, describir sus órbitas e identificar puntos fijos, puntos periódicos, conjuntos $\alpha$- y $\omega$-límite, y distinguir periodicidad, recurrencia y no errancia.
\item Comparar sistemas mediante conjugaciones topológicas, diferenciables y semiconjugaciones, y reconocer los invariantes preservados por cada tipo de equivalencia.
\item Describir la dinámica global de difeomorfismos del intervalo y aplicar resultados de linealización local cerca de puntos fijos contractivos o hiperbólicos.
\item Construir conjuntos de Cantor definidos dinámicamente, asociar itinerarios simbólicos a puntos y órbitas, y codificar mapas expansores por ramas completas mediante particiones generadoras.
\item Construir medidas invariantes para aplicaciones continuas de espacios compactos, aplicar el teorema de recurrencia de Poincaré y estudiar la convergencia de promedios temporales mediante el teorema ergódico de Birkhoff.
\item Probar la invariancia y ergodicidad de la medida de Lebesgue para mapas afines por ramas completas, analizar el mapa de Gauss y utilizar estimaciones de distorsión acotada para mapas expansores no lineales.
\item Analizar puntos fijos, órbitas periódicas y bifurcaciones en la familia logística, estudiar el papel de la órbita crítica y relacionar observaciones numéricas con propiedades topológicas y ergódicas.
\item Formular argumentos rigurosos, consultar bibliografía especializada y utilizar experimentos computacionales como apoyo para la exploración y la formulación de conjeturas.
\end{enumerate}

\subsubsection*{Contenidos}

\begin{enumerate}
\item \textbf{Sistemas dinámicos en tiempo discreto y continuo (2 semanas):}
\begin{enumerate}
    \item Sistemas dinámicos como acciones de $\mathbb{N}$, $\mathbb{Z}$ y $\mathbb{R}$; iteraciones y flujos.
    \item Órbitas positivas, negativas y completas; puntos fijos y de equilibrio; puntos y órbitas periódicas.
    \item Estabilidad y atracción; relación entre un flujo y sus mapas de tiempo fijo.
    \item Ejemplos elementales de ecuaciones diferenciales autónomas y mapas del intervalo.
\end{enumerate}

\item \textbf{Conjuntos límite, recurrencia y conjugaciones (2 semanas):}
\begin{enumerate}
    \item Conjuntos $\alpha$-límite y $\omega$-límite; puntos recurrentes y no errantes.
    \item Conjuntos invariantes y minimales.
    \item Conjugación topológica, semiconjugación y conjugación de flujos; invariantes topológicos.
    \item Conjugación diferenciable; multiplicadores de puntos periódicos y obstrucciones a la conjugación diferenciable.
\end{enumerate}

\item \textbf{Difeomorfismos del intervalo y linealización local (2 semanas):}
\begin{enumerate}
    \item Homeomorfismos y difeomorfismos de intervalos compactos; orientación y monotonía.
    \item Dinámica de difeomorfismos que preservan o invierten la orientación.
    \item Contracciones y teorema del punto fijo; puntos fijos hiperbólicos.
    \item Linealización topológica y diferenciable en dimensión uno; resultados de Hartman--Grobman en el contexto unidimensional.
\end{enumerate}

\item \textbf{Codificación simbólica (2 semanas):}
\begin{enumerate}
    \item Espacios de sucesiones y topología producto; desplazamientos unilateral y bilateral.
    \item Cilindros, puntos periódicos y órbitas densas; sistemas iterados de funciones.
    \item Construcción y codificación de conjuntos de Cantor definidos dinámicamente; ramas inversas e itinerarios.
    \item Mapas expansores por ramas completas; particiones generadoras y conjugación o semiconjugación con el desplazamiento.
\end{enumerate}

\item \textbf{Medidas invariantes y teoría ergódica (3 semanas):}
\begin{enumerate}
    \item Espacios de probabilidad y transformaciones medibles; medidas invariantes.
    \item Medidas de Dirac sobre órbitas periódicas; medidas empíricas.
    \item Construcción de Krylov--Bogolyubov para aplicaciones continuas de espacios métricos compactos.
    \item Teorema de recurrencia de Poincaré; ergodicidad, conjuntos invariantes y funciones invariantes.
    \item Teorema ergódico de Birkhoff.
\end{enumerate}

\item \textbf{Mapas por ramas, mapa de Gauss y distorsión acotada (2 semanas):}
\begin{enumerate}
    \item Transformaciones afines por ramas completas; invariancia y ergodicidad de la medida de Lebesgue.
    \item Operador de transferencia de Perron--Frobenius; el mapa de Gauss y fracciones continuas; medida de Gauss.
    \item Distorsión y cocientes de derivadas; lemas de distorsión acotada.
    \item Medidas invariantes absolutamente continuas para mapas expansores por ramas completas; ergodicidad respecto de la medida invariante y ergodicidad no singular respecto de Lebesgue.
\end{enumerate}

\item \textbf{La familia cuadrática (2 semanas):}
\begin{enumerate}
    \item Familia logística $q_a(x)=ax(1-x)$, con $0\leq a\leq4$, y relación con la familia $x\mapsto x^2+c$.
    \item Puntos fijos, estabilidad, órbitas periódicas y bifurcaciones; órbita del punto crítico.
    \item Intervalos invariantes, atractores, itinerarios y codificación simbólica; conjuntos de Cantor hiperbólicos.
    \item Ventanas periódicas, diagramas de bifurcación y promedios de Birkhoff.
\end{enumerate}
\end{enumerate}

\subsubsection*{Metodología}

\noindent\textbf{Lineamientos metodológicos:} La persona docente a cargo de este curso debe respetar las siguientes pautas:
\begin{enumerate}
    \item Combinar exposición teórica rigurosa con discusión de ejemplos concretos y resolución de problemas.
    \item Destacar las hipótesis de los teoremas, las ideas centrales de las demostraciones y la relación entre los enfoques topológico, diferenciable, simbólico y medible.
    \item Promover la comunicación clara de argumentos matemáticos en forma oral y escrita.
    \item Fomentar el trabajo independiente mediante listas de ejercicios y lecturas guiadas de la bibliografía.
\end{enumerate}

\textbf{Sugerencias metodológicas:} Adicionalmente, se recomienda:
\begin{enumerate}
    \item Desarrollar talleres en los que las personas estudiantes calculen órbitas, construyan conjugaciones, verifiquen invariancia de medidas y elaboren demostraciones breves.
    \item Utilizar Python, SageMath u otras herramientas disponibles para visualizar diagramas de telaraña, aproximar conjuntos de Cantor, representar itinerarios, construir diagramas de bifurcación y calcular promedios de Birkhoff; la experimentación computacional tendrá función exploratoria y no sustituirá la argumentación matemática.
    \item Solicitar que, durante el trabajo independiente, se revisen definiciones y demostraciones antes de cada clase, se resuelvan ejercicios de manera constante, se registren dudas para discutirlas en consulta o talleres, se lean previamente los fragmentos indicados de la bibliografía y se contrasten los experimentos numéricos con los resultados teóricos.
    \item Apoyarse en la plataforma institucional para distribuir materiales, listas de ejercicios, lecturas y recursos computacionales cuando corresponda.
\end{enumerate}

\nocite{BrSt02}
\nocite{HaKa03}
\nocite{KaHa95}
\nocite{Rob99}
\nocite{Sam05}
\nocite{SamUr}
\nocite{ViaOl16}
\nocite{Wal82}
\nocite{MeSt93}

\printcursosbibliography
\end{refsection}
